% Options for packages loaded elsewhere
\PassOptionsToPackage{unicode}{hyperref}
\PassOptionsToPackage{hyphens}{url}
\documentclass[
  ignorenonframetext,
]{beamer}
\newif\ifbibliography
\usepackage{pgfpages}
\setbeamertemplate{caption}[numbered]
\setbeamertemplate{caption label separator}{: }
\setbeamercolor{caption name}{fg=normal text.fg}
\beamertemplatenavigationsymbolsempty
% remove section numbering
\setbeamertemplate{part page}{
  \centering
  \begin{beamercolorbox}[sep=16pt,center]{part title}
    \usebeamerfont{part title}\insertpart\par
  \end{beamercolorbox}
}
\setbeamertemplate{section page}{
  \centering
  \begin{beamercolorbox}[sep=12pt,center]{section title}
    \usebeamerfont{section title}\insertsection\par
  \end{beamercolorbox}
}
\setbeamertemplate{subsection page}{
  \centering
  \begin{beamercolorbox}[sep=8pt,center]{subsection title}
    \usebeamerfont{subsection title}\insertsubsection\par
  \end{beamercolorbox}
}
% Prevent slide breaks in the middle of a paragraph
\widowpenalties 1 10000
\raggedbottom
\AtBeginPart{
  \frame{\partpage}
}
\AtBeginSection{
  \ifbibliography
  \else
    \frame{\sectionpage}
  \fi
}
\AtBeginSubsection{
  \frame{\subsectionpage}
}
\usepackage{iftex}
\ifPDFTeX
  \usepackage[T1]{fontenc}
  \usepackage[utf8]{inputenc}
  \usepackage{textcomp} % provide euro and other symbols
\else % if luatex or xetex
  \usepackage{unicode-math} % this also loads fontspec
  \defaultfontfeatures{Scale=MatchLowercase}
  \defaultfontfeatures[\rmfamily]{Ligatures=TeX,Scale=1}
\fi
\usepackage{lmodern}
\ifPDFTeX\else
  % xetex/luatex font selection
\fi
% Use upquote if available, for straight quotes in verbatim environments
\IfFileExists{upquote.sty}{\usepackage{upquote}}{}
\IfFileExists{microtype.sty}{% use microtype if available
  \usepackage[]{microtype}
  \UseMicrotypeSet[protrusion]{basicmath} % disable protrusion for tt fonts
}{}
\makeatletter
\@ifundefined{KOMAClassName}{% if non-KOMA class
  \IfFileExists{parskip.sty}{%
    \usepackage{parskip}
  }{% else
    \setlength{\parindent}{0pt}
    \setlength{\parskip}{6pt plus 2pt minus 1pt}}
}{% if KOMA class
  \KOMAoptions{parskip=half}}
\makeatother
\usepackage{longtable,booktabs,array}
\usepackage{caption}
\captionsetup[table]{skip=6pt}
\newcounter{none} % for unnumbered tables
\usepackage{calc} % for calculating minipage widths
\usepackage{caption}
% Make caption package work with longtable
\makeatletter
\def\fnum@table{\tablename~\thetable}
\makeatother
\setlength{\emergencystretch}{3em} % prevent overfull lines
\providecommand{\tightlist}{%
  \setlength{\itemsep}{0pt}\setlength{\parskip}{0pt}}
% Auto-generated by infrastructure/rendering/_pdf_latex_helpers.py::extract_math_font_preamble.
% Minimal header passed to Pandoc via -H so Beamer slides pick up the same math font
% as the combined-PDF path. Adjust the manuscript's preamble.md to change the font globally.
\usepackage{unicode-math}
\setmathfont{latinmodern-math.otf}

% Slide layout defaults for warning-clean scientific decks.
\usepackage{etoolbox}
\IfFileExists{xurl.sty}{\usepackage{xurl}}{}
\IfFileExists{seqsplit.sty}{\usepackage{seqsplit}}{\newcommand{\seqsplit}[1]{#1}}
\protected\def\breakseq#1{\seqsplit{#1}}
\protected\def\breaktt#1{\begingroup\ttfamily\seqsplit{#1}\endgroup}
\setlength{\emergencystretch}{6em}
\tolerance=5000
\hbadness=10000
\hfuzz=1pt
\setlength{\tabcolsep}{2pt}
\AtBeginEnvironment{longtable}{\tiny\renewcommand{\arraystretch}{0.86}\setlength{\tabcolsep}{1pt}}
\AtBeginEnvironment{tabular}{\tiny\renewcommand{\arraystretch}{0.86}\setlength{\tabcolsep}{1pt}}
\AtBeginEnvironment{equation}{\tiny}
\AtBeginEnvironment{equation*}{\tiny}
\AtBeginEnvironment{align}{\tiny}
\AtBeginEnvironment{align*}{\tiny}
\AtBeginEnvironment{itemize}{\footnotesize}
\AtBeginEnvironment{enumerate}{\footnotesize}
\AtBeginEnvironment{description}{\footnotesize}
\setbeamerfont{caption}{size=\tiny}
\setbeamerfont{caption name}{size=\tiny}
\setbeamerfont{normal text}{size=\small}
\setbeamerfont{frametitle}{size=\small}
\setbeamerfont{section title}{size=\footnotesize}
\setbeamerfont{subsection title}{size=\footnotesize}
\setbeamertemplate{section page}{%
  \centering
  \begin{beamercolorbox}[sep=12pt,center,wd=\paperwidth]{section title}
    \parbox{0.86\paperwidth}{\centering\usebeamerfont{section title}\insertsection\par}
  \end{beamercolorbox}
}
\setbeamertemplate{subsection page}{%
  \centering
  \begin{beamercolorbox}[sep=8pt,center,wd=\paperwidth]{subsection title}
    \parbox{0.86\paperwidth}{\centering\usebeamerfont{subsection title}\insertsubsection\par}
  \end{beamercolorbox}
}
\setlength{\abovecaptionskip}{2pt}
\setlength{\belowcaptionskip}{0pt}

% Natbib and cross-reference fallbacks — slides don't load natbib
% or cleveref, but combined-PDF manuscript prose may emit these
% commands. The fallback renders citations as a bracketed key list
% and cross-references as detokenized labels so slides don't fail on
% undefined control sequences or raw underscores. \providecommand is
% a no-op if packages load later.
\providecommand{\citep}[1]{[#1]}
\providecommand{\citet}[1]{#1}
\providecommand{\citealp}[1]{#1}
\providecommand{\citeauthor}[1]{#1}
\providecommand{\citeyear}[1]{#1}
\providecommand{\cref}[1]{\texttt{\detokenize{#1}}}
\providecommand{\Cref}[1]{\texttt{\detokenize{#1}}}

\newtheorem{proposition}{Proposition}
\newtheorem{hypothesis}{Hypothesis}
\usepackage{bookmark}
\IfFileExists{xurl.sty}{\usepackage{xurl}}{} % add URL line breaks if available
\urlstyle{same}
% fallback for those not using the hyperref driver hyperxmp:
\makeatletter
\@ifundefined{xmpquote}{\newcommand{\xmpquote}[1]{#1}}{}
\makeatother
\hypersetup{
  hidelinks,
  pdfcreator={LaTeX via pandoc}}

\author{\texorpdfstring{}{}}
\date{}

\begin{document}

\section{Related Work: SkillTree Delivery, Security Posture, Reporting,
and Open
Science}\label{sec:related-work-skilltree-security-reporting-open-science}

\begin{frame}[allowframebreaks]{SkillTree delivery as prerequisite
learning and gamified practice context}
\protect\phantomsection\label{skilltree-delivery-as-prerequisite-learning-and-gamified-practice-context}
SkillTree supplies the delivery substrate: Project, Subject, Skill,
learning-path dependencies, quizzes, and report endpoints (National
Security Agency 2020, 2026b; SkillTree 2026). Active Skillference uses
those concepts directly rather than inventing a new learning platform.
The graph layer adds prerequisite validation and connectivity checks
because the learning object is the ordered dependency structure, not
only the text attached to each node (Bloom 1956; Bloom 1968; Liang et
al. 2018; Doignon and Falmagne 1985; Falmagne and Doignon 2011).

\end{frame}

\begin{frame}[allowframebreaks]{SkillTree delivery as prerequisite
learning and gamified practice context}
Gamification is relevant but not itself the contribution; SkillTree
provides the gamified runtime, and the project contributes the verified
content and graph (Dicheva et al. 2015). Visual skill-tree and gamified
microlearning sources provide useful delivery precedent for making
dependencies and progress inspectable, but they do not constitute
learner-outcome evidence for this export (Burau 2022; Septiani and
Rosmansyah 2021). Constructive alignment and formative-assessment
sources supply the education-facing language for this design:
objectives, activities, feedback, and quiz gates are treated as aligned
checkpoints, not as outcome evidence (Biggs 1996; Bloom 1968; Black and
Wiliam 1998; Shute 2008).
\end{frame}

\begin{frame}[fragile,allowframebreaks]{SkillTree provenance, deployment
security, and access-control boundaries}
\protect\phantomsection\label{skilltree-provenance-deployment-security-and-access-control-boundaries}
SkillTree's origin matters for the delivery argument. NSA describes
SkillTree as an internally developed open-source training platform, and
the public \texttt{skills-service} repository and NSA open-source portal
identify the service as public government-origin software rather than a
generic third-party widget (National Security Agency 2020, 2026b,
2026a). Active Skillference therefore treats SkillTree as a real
delivery substrate with explicit platform ownership boundaries:
SkillTree owns learner progress, quiz state, authentication, ranking and
reporting semantics, while this project owns the curriculum graph,
provenance gates, manuscript evidence and exported SkillTree payload.

\end{frame}

\begin{frame}[fragile,allowframebreaks]{SkillTree provenance, deployment
security, and access-control boundaries}
That provenance raises the bar for deployment language. A
national-security posture would still require organization-specific
threat modeling, zero-trust access decisions, supply-chain provenance,
software-development process, artifact signing, SBOM visibility, and
adversary-technique coverage beyond this educational repository (Rose et
al. 2020; National Institute of Standards and Technology 2022; SLSA
Framework 2026; Sigstore 2026; Cybersecurity and Infrastructure Security
Agency 2026; MITRE 2026). Active Skillference can support defensible
training review because claims, figures, source roles, dashboard
artifacts and SkillTree handoff routes are inspectable; it does not
certify operator competence, harden a production deployment, replace
APT-aware threat modeling, or imply NSA approval of the curriculum.
\end{frame}

\begin{frame}[fragile,allowframebreaks]{Reproducible reporting,
accessible artifacts, and source-role auditing}
\protect\phantomsection\label{reproducible-reporting-accessible-artifacts-and-source-role-auditing}
The content-provenance binding is closest in spirit to literate
programming and literate statistical reporting: code and explanation are
kept mutually inspectable, and R/knitr- style workflows bind prose to
executable analysis so reported numbers can be regenerated (Knuth 1984;
Xie 2015). Peng's reproducible-research argument is the reporting
boundary used here: when a claim is computational, readers should be
able to rerun the workflow that produced it (Peng 2011; Sandve et al.
2013). Scientific-computing practice adds the software-maintenance side
of the same argument: versioned code, automated checks, and tested
workflows are part of the scholarly surface when the artifact is
executable (Wilson et al. 2014). W3C PROV gives the formal provenance
vocabulary for entities, activities, agents, derivations, bundles, and
\end{frame}

\begin{frame}[fragile,allowframebreaks]{Reproducible reporting,
accessible artifacts, and source-role auditing}
validity constraints, while reproducible-science sources keep the claim
boundary tied to transparent methods rather than trust language (Groth
and Moreau 2013; Moreau and Missier 2013; Cheney et al. 2013; Munafò et
al. 2017). FAIR and transparency guidelines extend that boundary from
rerunnable computation to metadata and open-science claim discipline
(Wilkinson et al. 2016; Nosek et al. 2015). Active Skillference differs
by keeping a central claim registry, applying the same gate to
curriculum content and manuscript prose, and refusing export when any
token is unbacked. The DAG engine uses standard graph algorithms for
cycle detection and topological ordering (Cormen et al. 2022). Figures
use an accessible categorical palette rather than a single-hue theme
(Wong 2011). The figure-registry validator also treats alt text and
semantic caption metadata as reporting obligations, aligning the
\end{frame}

\begin{frame}[fragile,allowframebreaks]{Reproducible reporting,
accessible artifacts, and source-role auditing}
artifact layer with ACM/SIGCHI, Section 508 and WCAG guidance that
accessible figures and web content should carry meaningful alternatives
rather than empty labels (ACM DIS 2023 Accessibility Chairs 2023;
Section508.gov 2026; World Wide Web Consortium 2023).

Release metadata and license notices are treated as delivery evidence,
not as scientific or learner-efficacy support. Software citation
principles, Citation File Format and CodeMeta guide the citation and
metadata surfaces (Smith et al. 2016; Citation File Format Project 2026;
CodeMeta Project 2026). GitHub and Zenodo documentation bound the
repository and DOI workflow to an actual public target and deposit
(GitHub Docs 2026; Zenodo 2026). SPDX, OSI and REUSE license-notice
references support the MIT and third-party notice signposting (SPDX
Workgroup 2026; Open Source Initiative 2026; Free Software Foundation
Europe 2024). These references do not provide legal advice, do not
\end{frame}

\begin{frame}[fragile,allowframebreaks]{Reproducible reporting,
accessible artifacts, and source-role auditing}
relicense third-party course materials or vendored browser bundles, and
do not turn a static review checkout into a production deployment.

The scholarship audit extends that reproducible-reporting idea from
numbers to source coverage. Ordinary citation validators can show that
\texttt{\breakseq{[@key]}} resolves, but not that the key has a stated role in
the argument. The source matrix makes that role explicit and keeps the
bibliography aligned with the manuscript's actual lanes. The frontier
gap report adds a narrower check for newly emphasized source areas: a
pedagogy, assessment, formal provenance, workflow provenance,
knowledge-space, cognitive-diagnosis, message-passing,
predictive-coding, design-based evaluation, open-science, security,
accessibility, reproducibility, or course-material source must be cited
and tied to an explicit claim-support class and, where needed, a
boundary row before it can count as a scholarship upgrade.
\end{frame}

\begin{frame}[allowframebreaks]{Open-science education ecosystem,
interoperability, and community handoff}
\protect\phantomsection\label{open-science-education-ecosystem-interoperability-and-community-handoff}
Active Skillference also sits inside a broader open-science educational
ecosystem. The Active Inference Institute publishes a public resource
hub with community pathways, verified public repositories, and a
documented projection boundary; it also publishes public course
infrastructure, a START project for tailored Active Inference research
and training curricula, and an institute website that frames education
as part of an open-science mission (Active Inference Institute 2026b,
2026e, 2026a, 2026c).

\end{frame}

\begin{frame}[allowframebreaks]{Open-science education ecosystem,
interoperability, and community handoff}
The cognitive knowledge-base repository is treated as course-material
provenance for generated Learn pointers, not as copied prose or as an
endorsement claim (Active Inference Institute 2026d). This project does
not replace those efforts; it contributes a provenance-bound SkillTree
export, quiz gate, and verifier suite that can interoperate with public
learning paths. Recent educational Active Inference scholarship also
motivates prepared learning environments and active exploration as
topics for future classroom evaluation, while this release limits itself
to the inspectable curriculum and delivery contract (Di Paolo et al.
2024).
\end{frame}

\end{document}
